\documentclass[11pt,a4paper,oneside]{article}

\usepackage[T1]{fontenc}
\usepackage[utf8]{inputenc}
\usepackage[spanish,es-noquoting,es-nodecimaldot]{babel}
\usepackage{lmodern}
\usepackage{microtype}
\usepackage[a4paper,margin=1.18in,headheight=15pt]{geometry}
\usepackage{amsmath,amssymb,amsfonts,mathtools}
\usepackage{enumitem}
\usepackage{xcolor}
\usepackage{hyperref}
\usepackage{setspace}
\usepackage{fancyhdr}
\usepackage{titlesec}
\usepackage[most]{tcolorbox}

\hypersetup{
  colorlinks=true,
  linkcolor=blue!50!black,
  urlcolor=teal!60!black,
  pdftitle={Tarea: Ecuaciones Diferenciales Lineales de Primer Orden},
  pdfauthor={César Galindo}
}

\theoremstyle{definition}
\newtheorem{tarea}{Ejercicio}

\definecolor{econblue}{RGB}{25,62,114}
\definecolor{econbrown}{RGB}{138,90,47}
\pagestyle{fancy}
\fancyhf{}
\fancyhead[L]{\textsc{Cálculo Integral -- Facultad de Empresariales}}
\fancyhead[R]{\textsc{César Galindo}}
\fancyfoot[C]{\thepage}
\renewcommand{\headrulewidth}{0.35pt}
\renewcommand{\footrulewidth}{0pt}
\titleformat{\section}{\Large\bfseries\color{econblue}}{\thesection.}{0.65em}{}
\titlespacing*{\section}{0pt}{1.9em}{1.05em}
\tcbset{colback=white,colframe=econblue!65!black,boxrule=0.5pt,arc=1.5pt}
\setlength{\parskip}{0.85em}
\setstretch{1.13}
\setlength{\parindent}{0pt}

\newtcolorbox{ideaBox}[1][]{enhanced,breakable,colback=blue!2,colframe=econblue,
  borderline west={1.5pt}{0pt}{econblue},title=Instrucciones,fonttitle=\bfseries,#1}

\title{\vspace{-1.5cm}\textbf{\Large Tarea}\\[2pt]
\large Ecuaciones Diferenciales Lineales de Primer Orden}
\author{César Galindo\\ \small Cálculo Integral \textperiodcentered\ Facultad de Empresariales}
\date{\small Otoño 2026}

\begin{document}
\maketitle
\vspace{-0.8cm}

\begin{ideaBox}
Resolver los cinco ejercicios siguientes, mostrando todos los pasos
intermedios. Apóyense en las fórmulas y teoremas de la nota de clase
\emph{Ecuaciones Diferenciales Lineales de Primer Orden} (factor
integrante, ecuación de Bernoulli, y las tres aplicaciones financieras
de la Sección de Aplicaciones): en particular, en los Ejercicios~1--3 de
esa nota, ya resueltos paso a paso, como modelo de la presentación
esperada.

\smallskip
\textbf{Fecha de entrega:} \underline{\hspace{4.5cm}}
\end{ideaBox}

\begin{tarea}[Factor integrante con forzamiento oscilatorio]
Resolver $y' + y = 2\cos x$ con $y(0)=0$.
\emph{Sugerencia:} $\displaystyle\int e^x\cos x\,dx =
\tfrac12 e^x(\cos x+\sin x)+C$; verifiquen la fórmula derivando antes de
usarla.
\end{tarea}

\begin{tarea}[Bernoulli con $n=2$]
Resolver $y' - y = xy^2$.
\emph{Sugerencia:} usen la sustitución $v=y^{-1}$; la ecuación lineal
resultante para $v$ se resuelve con un factor integrante exponencial.
\end{tarea}

\begin{tarea}[Aplicación financiera: fondo de amortización de un bono corporativo]
Una empresa debe pagar \$5{,}000{,}000 al vencimiento de un bono a 10
años, y constituye hoy un \emph{sinking fund} con $A_0=0$ que gana un
rendimiento continuo $r=4.5\%$ anual, alimentado con aportaciones
continuas a tasa constante $c$.
\begin{enumerate}[label=(\alph*)]
  \item Calcular la aportación anual continua $c$ necesaria para que el
  fondo alcance exactamente \$5{,}000{,}000 en $t=10$.
  \item Si la empresa parte en cambio con $A_0=\$800{,}000$ ya ahorrados
  y mantiene la misma $c$ del inciso (a), calcular cuánto tendrá el
  fondo en $t=10$, y por cuánto excede la meta.
\end{enumerate}
\end{tarea}

\begin{tarea}[Aplicación financiera: dotación perpetua para una cátedra universitaria]
Un donante quiere financiar a perpetuidad una cátedra universitaria que
cuesta \$300{,}000 continuos por año, con un fondo invertido a un
rendimiento requerido continuo $r=5\%$.
\begin{enumerate}[label=(\alph*)]
  \item Calcular el monto $V$ que el donante debe aportar hoy para que
  el fondo pague la cátedra \textbf{para siempre} sin agotarse.
  \item Si el donante sólo dispone de $B_0=\$4{,}500{,}000$ y aun así se
  retiran continuamente \$300{,}000 por año, mostrar que el fondo se
  agota en tiempo finito y calcular en cuántos años.
\end{enumerate}
\end{tarea}

\begin{tarea}[Aplicación financiera: reversión del diferencial de crédito]
El diferencial de crédito de un bono corporativo está hoy en
$r_0=3.5\%$ sobre la tasa libre de riesgo. El mercado espera que
revierta a su nivel de largo plazo $b=1.8\%$ a una velocidad $a=0.4$
(por año).
\begin{enumerate}[label=(\alph*)]
  \item Calcular el diferencial proyectado a 1 y a 3 años.
  \item ¿Cuántos años se necesitan para recorrer el 80\% de la
  distancia entre $r_0$ y $b$?
\end{enumerate}
\end{tarea}

\end{document}
